\section{Supporting Information}
Here we include some peripheral elements of work that support
 the paper.

\subsection{Notes on Measurement Error Propagation}
\label{subsec:NoMEP}
\textcolor{ForestGreen}{
To evaluate uncertainty associated with the percentage mass loss
 metric we undertook some error propagation analysis by modelling
 the process of obtaining values for this metric.
This modelling was carried out in the Python programming language
 and is detailed in a Jupyter notebook held within the Github
 repository associated with this paper, it can be found
 under the `Sandbox/UncertaintyProgagation' directory and is
 entitled `UncertaintyModel.ipynb'.
}

\textcolor{ForestGreen}{
By setting the initial mass of a given sample at 3~\unit{\gram}
 and allowing the final mass to vary in accordance with a uniform
 distribution between a lower and upper bound (as deduced by
 from the two samples evaluated exhibiting the maximal and minimal
 mass losses across all trials) we could simulate the drawing of
 a random sample from across the parameter space each having its own
 unique mass loss.
Additionally, to model the measurement error which manifests itself
 during the process of evaluating the initial mass and final mass
 values we could employ normal distributions with standard deviations
 equal to the repeatability of the scales employed (0.01~\unit{\gram}).
By coupling these concepts together and undertaking a Monte Carlo
 sampling exercise we could obtain estimates of the average
 uncertainty associated with any given percentage mass loss value
 considered in the present work; this was equal to 
 0.44~\% points (1\(\sigma\)=0.01\% points).
}


\subsection{Notes on Process-Energy Calculations}
\label{subsec:NoPEC}

\textcolor{blue}{
Although this paper has focussed upon branching polyester
 thermosets as an alternative to unsaturated polyester resin
 derived thermosets on account of the relative ease with
 which these materials may be prepared from bio-based feedstocks,
 we acknowledge the fact that this does not inherently make this
 family of polymer more sustainable than their conventional
 counterparts.
To prove this, one would first have to choose a suitable 
 quantitative definition of sustainability (e.g. embodied carbon),
 before engaging in a comprehensive life-cycle analysis for both
 the methods by which thermosets were manufactured.
}

\textcolor{blue}{
Although by no means a full life-cycle analysis of either route,
 we have undertaken some modelling of the two system with a view
 to elucidating whether the prolonged thermal cure time used for
 branched poly(glycerol citrate itaconate) thermosets is likely
 to impose an energy penalty relative to the unsaturated polyester 
 resin route.
This model is presented within a Jupyter notebook file 
 entitled `calcs\_restructured\_uncertainty.ipynb' (within the 
 `Sandbox/EnergyCalculations' directory) in
 the github repository associated with this paper.
}

\textcolor{blue}{
Here, we establish two separate models, one for the branched system
 and the other for the unsaturated polyester system.
These models both consider the energy required to form prepolymers
 from feedstocks and the subsequent energy required to cure these
 prepolymers.
Each model is based on a set of parameters, each parameter being
 described by a triangular distribution, which allows propagation 
 of uncertainty to be considered over the course of a large number
 of simulations in a Monte Carlo fashion.
}

\textcolor{blue}{
On conclusion of 20,000 Monte Carlo draws, it was concluded that
 the more energy intensive route was that of branched polyester
 thermoset manufacture (requiring an average of 5.4x the 
 kWh~kg\textsuperscript{-1} of thermoset).
Table~\ref{tab:SummaryOfEnergyModelling} provides an overview of
 this exercise, and shows how the progagated error bounds 
 surrounding the median values.
Estimates of kWh~kg\textsuperscript{-1} yielded by the branched 
 model were particularly sensitive toward power required to
 maintain the oven temperature as well as the cure duration,
 whilst the unsaturated polyester model was most sensitive toward
 energy required to maintain reactor temperature, auxiliary power
 requirements, and heat capacity of various elements.
 }

\small
\begin{figure*}[htb]
    \begin{center}
        \begin{tblr}{
            row{1} = {c},
            cell{2}{2} = {r},
            cell{2}{3} = {r},
            cell{2}{4} = {r},
            cell{3}{2} = {r},
            cell{3}{3} = {r},
            cell{3}{4} = {r},
            cell{4}{2} = {r},
            cell{4}{3} = {r},
            cell{4}{4} = {r},
            vline{2} = {-}{},
            hline{2} = {-}{},
            hline{4} = {-}{dashed},
            }
            Metric                                                  & 5\textsuperscript{th} Percentile & Median & 95\textsuperscript{th} Percentile \\
            Branched Polyester Route (kWh kg\textsuperscript{-1})   & 24.4                             & 49.9   & 84.1           \\
            Conventional UPR Route (kWh kg\textsuperscript{-1})     & 5.4                              & 9.2    & 14.6            \\
            UPR/Branched Ratio                                      & 0.087                            & 0.186  & 0.431           
        \end{tblr}
        \captionof{table}{
                \footnotesize{
                    Summary of process-energy 
                    modelling exercise.
                }
            }
        \label{tab:SummaryOfEnergyModelling}
    \end{center}
\end{figure*}
\normalsize